
%----- Math stuff -------------------------------------------------------
\usepackage{mathcomp}
\usepackage{icomma}
\usepackage{amsmath,amssymb,amsfonts,bbm}
\usepackage{amsthm}   % AMS-Theoremumgebung
\usepackage{stmaryrd}	% Symbole, u. A. \lightning
\usepackage{units}
\usepackage{sistyle} % für physikalische Einheiten z. B. \SI{1000}{m/s^2}

%----- Aufzählung, Reihe ------------------------------------------------
\newcommand{\folge}[2][k]{\mbox{\ensuremath{#2_1, \ldots, #2_{#1}}}}
%------------------------------------------------------------------------

%----- Zusätzliche Mathematische Funktionen -----------------------------
\newcommand*\abs[1]{\left|#1\right|} % abs{-x}
\newcommand*\ueber[2]{\left({#1\atop#2}\right)} % \ueber{n}{k}
%------------------------------------------------------------------------

%----- Befehl \textfrac{z}{n} - für schräge Brüche in Textpassagen ------
\newcommand*{\textfrac}[2]{\ensuremath{\kern.1em
\raise.5ex\hbox{\the\scriptfont0 #1}\kern-.1em
/\kern-.15em\lower.25ex\hbox{\the\scriptfont0 #2}}}
%------------------------------------------------------------------------

%----- efrag - für Brüche mit (aufrecht gesetzten) Einheiten ------------
\newcommand{\efrac}[2]{\frac{\mathrm{#1}}{\mathrm{#2}}}
%------------------------------------------------------------------------

%----- Mathematische Symbole --------------------------------------------
% Entspricht-Symbol
%\newcommand{\entspricht}{\mathrel{\widehat{=}}}
\newcommand{\entspricht}{\stackrel{\scriptscriptstyle\wedge}{=}}
% Langer Pfeil rechts mit optionalem Argument oben
\newcommand{\rto}[1][~]{\ensuremath{\xrightarrow{\;#1\;}}}
% Äquivalenzpfeile
\newcommand{\same}{\ensuremath{\Leftrightarrow}}
\newcommand{\qsame}{\quad\same\quad}
\newcommand{\lsame}{\ensuremath{\Longleftrightarrow}}
\newcommand{\qlsame}{\quad\lsame\quad}
\newcommand{\dsame}{\ensuremath{:\Leftrightarrow}}
\newcommand{\qdsame}{\quad\dsame\quad}
\newcommand{\ldsame}{\ensuremath{:\Longleftrightarrow}}
\newcommand{\qldsame}{\quad\ldsame\quad}
\newcommand{\samedef}{\ensuremath{\us {\text{def}} \Leftrightarrow}}
\newcommand{\qsamedef}{\quad\samedef\quad}
\newcommand{\lsamedef}{\ensuremath{\us {\text{def}} \Longleftrightarrow}}
\newcommand{\qlsamedef}{\quad\lsamedef\quad}
\newcommand{\lto}{\longrightarrow}
% Folgepfeil
\newcommand{\qTo}{\quad\To\quad}
\newcommand{\sTo}{\ \To\ }
\newcommand{\To}{\ensuremath{\Rightarrow}}
% Logisches Und und Oder
\newcommand{\sand}{\ensuremath{\; \land \;}}
\newcommand{\sor}{\ensuremath{\; \lor \;}}
\renewcommand{\*}{\ensuremath{\cdot}}
\newcommand{\od}{\vee}
\newcommand{\n}{\neg}
\newcommand{\un}{\wedge}
% einige Negationen
\newcommand{\nin}{\not\in}
% Bei Integral, Summe,... die Grenzen über bzw. unter das Zeichen
\newcommand{\I}{\int\limits}
\newcommand{\Sum}{\sum\limits}
\newcommand{\Prod}{\prod\limits}
\newcommand{\Min}{\min\limits}
\newcommand{\Max}{\max\limits}
\newcommand{\Lim}{\lim\limits}
\newcommand{\Inf}{\inf\limits}
\newcommand{\Sup}{\sup\limits}
\newcommand{\Liminf}{\liminf\limits}
\newcommand{\Limsup}{\limsup\limits}
\renewcommand{\Cup}{\bigcup\limits}
\renewcommand{\Cap}{\bigcap\limits}
\newcommand{\Otimes}{\bigotimes\limits}
\newcommand{\Oplus}{\bigoplus\limits}
\newcommand{\isqm}{\ensuremath{\stackrel{\mathrm{?}}{=}}}
\renewcommand{\Re}{\text{Re}}
\renewcommand{\Im}{\text{Im}}
\newcommand{\const}{\text{const.}}
\newcommand{\Epot}{E_{\text{pot}}}
\newcommand{\Ekin}{E_{\text{kin}}}
%------------------------------------------------------------------------

%----- Logik-Symbole ----------------------------------------------------
\newcommand{\band}{\ensuremath{\wedge}}
\newcommand{\bAnd}{\ensuremath{\bigwedge\limits}}
\newcommand{\bor}{\ensuremath{\vee}}
\newcommand{\bOr}{\ensuremath{\bigvee\limits}}
\newcommand{\xor}{\ensuremath{\oplus}}
\newcommand{\Xor}{\ensuremath{\bigoplus\limits}}
%------------------------------------------------------------------------

%----- Mengensymbole ----------------------------------------------------
\newcommand{\N}{\ensuremath{\mathbbm{N}}}  % natürliche Zahlen
\newcommand{\Z}{\ensuremath{\mathbbm{Z}}}  % ganze Zahlen
\newcommand{\Q}{\ensuremath{\mathbbm{Q}}}  % rationale Zahlen
\newcommand{\R}{\ensuremath{\mathbbm{R}}}  % reelle Zahlen
\newcommand{\C}{\ensuremath{\mathbbm{C}}}  % komplexe Zahlen
\newcommand{\qH}{\ensuremath{\mathbbm{H}}} % Quaternionen
\newcommand{\B}{\ensuremath{\mathbbm{B}}}  % Binärzahlen
%------------------------------------------------------------------------

%----- alternative griechische Kleinbuchstaben verwenden ----------------
\newcommand*\defaultepsilon{\epsilon}
\newcommand*\defaulttheta{\theta}
\newcommand*\defaultphi{\phi}
\renewcommand*\epsilon{\varepsilon}
\renewcommand*\theta{\vartheta}
\renewcommand*\phi{\varphi}
\newcommand*\lamda{\lambda}
%------------------------------------------------------------------------

%----- Caligraphische Zeichen und Frakturbuchstaben ---------------------
\newcommand{\cA}{\mathcal{A}}
\newcommand{\cB}{\mathcal{B}}
\newcommand{\cC}{\mathcal{C}}
\newcommand{\cD}{\mathcal{D}}
\newcommand{\cE}{\mathcal{E}}
\newcommand{\cF}{\mathcal{F}}
\newcommand{\cG}{\mathcal{G}}
\newcommand{\cH}{\mathcal{H}}
\newcommand{\cI}{\mathcal{I}}
\newcommand{\cJ}{\mathcal{J}}
\newcommand{\cK}{\mathcal{K}}
\newcommand{\cL}{\mathcal{L}}
\newcommand{\cM}{\mathcal{M}}
\newcommand{\cN}{\mathcal{N}}
\newcommand{\cO}{\mathcal{O}}
\newcommand{\cP}{\mathcal{P}\,}
\newcommand{\cQ}{\mathcal{Q}}
\newcommand{\cR}{\mathcal{R}}
\newcommand{\cS}{\mathcal{S}}
\newcommand{\cT}{\mathcal{T}}
\newcommand{\cU}{\mathcal{U}}
\newcommand{\cV}{\mathcal{V}}
\newcommand{\cW}{\mathcal{W}}
\newcommand{\cX}{\mathcal{X}}
\newcommand{\cY}{\mathcal{Y}}
\newcommand{\cZ}{\mathcal{Z}}
% Fraktur-Buchstaben
\newcommand{\fA}{\ensuremath{\mathfrak A}}
\newcommand{\fB}{\ensuremath{\mathfrak B}}
\newcommand{\fC}{\ensuremath{\mathfrak C}}
\newcommand{\fD}{\ensuremath{\mathfrak D}}
\newcommand{\fE}{\ensuremath{\mathfrak E}}
\newcommand{\fF}{\ensuremath{\mathfrak F}}
\newcommand{\fG}{\ensuremath{\mathfrak G}}
\newcommand{\fH}{\ensuremath{\mathfrak H}}
\newcommand{\fI}{\ensuremath{\mathfrak I}}
\newcommand{\fJ}{\ensuremath{\mathfrak J}}
\newcommand{\fK}{\ensuremath{\mathfrak K}}
\newcommand{\fL}{\ensuremath{\mathfrak L}}
\newcommand{\fM}{\ensuremath{\mathfrak M}}
\newcommand{\fN}{\ensuremath{\mathfrak N}}
\newcommand{\fO}{\ensuremath{\mathfrak O}}
\newcommand{\fP}{\ensuremath{\mathfrak P}}
\newcommand{\fQ}{\ensuremath{\mathfrak Q}}
\newcommand{\fR}{\ensuremath{\mathfrak R}}
\newcommand{\fS}{\ensuremath{\mathfrak S}}
\newcommand{\fT}{\ensuremath{\mathfrak T}}
\newcommand{\fU}{\ensuremath{\mathfrak U}}
\newcommand{\fV}{\ensuremath{\mathfrak V}}
\newcommand{\fW}{\ensuremath{\mathfrak W}}
\newcommand{\fX}{\ensuremath{\mathfrak X}}
\newcommand{\fY}{\ensuremath{\mathfrak Y}}
\newcommand{\fZ}{\ensuremath{\mathfrak Z}}
%------------------------------------------------------------------------

%----- Arcus Cotangens, Logarithmus -------------------------------------
\newcommand{\arsin}{\operatorname{arsin}}
\newcommand{\arcos}{\operatorname{arcos}}
\newcommand{\arsinh}{\operatorname{arsinh}}
\newcommand{\arcosh}{\operatorname{arcosh}}
\newcommand{\artanh}{\operatorname{arctanh}}
\newcommand{\arcot}{\operatorname{arcot}}
\newcommand{\arcoth}{\operatorname{arcoth}}
\newcommand{\Log}{\operatorname{Log}}
\renewcommand{\exp}{\operatorname{e}}
%------------------------------------------------------------------------

%----- Bezeichnungen ----------------------------------------------------
\newcommand{\Mat}{\operatorname{Mat}} % Matrizen
\newcommand{\Hom}{\operatorname{Hom}} % Homomorphismus
\newcommand{\grad}{\operatorname{grad}} % Grad
\newcommand{\Grad}{\operatorname{Grad}} % Grad eines Polynoms
\newcommand{\Rang}{\operatorname{Rang}} % Rang einer Matrix
\newcommand{\Pot}{\operatorname{Pot}}   % Potenzmenge
\newcommand{\id}{\operatorname{id}}     % Identit"at
\newcommand{\ord}{\operatorname{ord}}   % Ordnung eines Elements
\newcommand{\Def}{\operatorname{Def}}   % Definitionsbereich
\newcommand{\Bild}{\operatorname{Bild}} % Bildbereich
\newcommand{\Feld}{\operatorname{Feld}} % Feld von einer Relation
\newcommand{\Kern}{\operatorname{Kern}} % Kern eines Homomorphismes
\newcommand{\Det}{\operatorname{Det}}   % Determinante
\newcommand{\Rad}{\operatorname{Rad}}   % Radikal
\newcommand{\kgV}{\operatorname{kgV}}   % Kleinster gemeinsame Vielfache
\newcommand{\ggT}{\operatorname{ggT}}   % Größter gemeinsame Teiler
\newcommand{\Spur}{\operatorname{Spur}} % Spur
\newcommand{\Char}{\operatorname{char}} % Charakteristik
\newcommand{\diverg}{\operatorname{div}}% Divergenz
\newcommand{\rot}{\operatorname{rot}}   % Rotation
\newcommand*{\diff}[1]{\mathrm{d}#1}    % Differentialoperator
%------------------------------------------------------------------------

%----- häufige Abkürzungen ----------------------------------------------
\newcommand{\oBdA}{o.\,B.\,d.\,A.\ }
\newcommand{\zB}{z.\,B.\ }
\newcommand{\zT}{z.\,T.\ }
\newcommand{\insbes}{insbes.\ }
\renewcommand{\dh}{d.\,h.\ }
%------------------------------------------------------------------------

%----- Physikalische Einheiten ------------------------------------------
\newcommand*{\micro}{\ensuremath{\mbox{\textmu}}}
\newcommand*{\mikro}{\ensuremath{\mbox{\textmu}}}
\newcommand*{\microsec}{\ensureupmath{\micro s}}
\newcommand{\Ohm}{\mathrm{\Omega}}
\newcommand{\Celsius}{\degC}
\SIdecimalsign{,}
\SIthousandsep{\,}
\SIproductsign{\cdot}
\SIgroupfourtrue
\SIobeyboldtrue
% siehe Paketdokumentation zu SIstyle
%------------------------------------------------------------------------

%----- Astronomische Symbole --------------------------------------------
\usepackage{wasysym}
% siehe "The comprehensive LaTeX Symbol List"
%------------------------------------------------------------------------

%----- Boolean expressions ----------------------------------------------
\usepackage{boolexpr}
